%!TEX root = ../../thesis.tex

\subsection{open62541}\label{sec:target-open62541}

% The stateful real-world target: an OPC~UA server stack, driven by protocol
% messages. This is where the "stateful" claim in the title is actually cashed in.

% TODO: Describe the target and the harness:
%   - the server binary that is rehosted, and the attack surface reached through it,
%   - the session protocol the agent must navigate --- HELLO, OpenSecureChannel,
%     CreateSession, ActivateSession and the various session and channel requests ---
%     and why reaching an activated session is a meaningful depth objective:
%     everything interesting in the protocol lies behind it,
%   - actions are message templates drawn from the target's corpus, so the agent
%     learns *which* message to send *when*; the environment rewrites the
%     session-coupling identifiers from the server's own responses. Be explicit that
%     this is a deliberate scoping decision: the agent is not learning message bytes,
%     and therefore this target evaluates protocol-state navigation, not input
%     synthesis,
%   - reward: a living reward for each step the server survives and a bonus on
%     session activation, which also ends the episode.
%
% This is the target that connects to the stateful-fuzzing literature of
% \cref{sec:fuzzing-stateful} and to the state-selection bandits of
% \cref{sec:rlfuzz-prior-work}, so name the comparable systems (AFLNet, SGFuzz) and
% say honestly what is and is not comparable about them --- they mutate message
% contents, which this formulation does not.
%
% Also state the scaling gap it exposes: a real stack has far more basic blocks than
% the hashed bucket array of \cref{sec:platform-coverage} can separate.
% Carry that into \cref{sec:discussion-threats}.
